\chapter{Procesamiento del Lenguaje Natural Cuántico (\glsentrytext{qnlp})}

\section{Motivación para el uso de \glsentrytext{qnlp}}

Tal y como se expuso en el capítulo precedente, la implementación del \glsentrytext{nlp} mediante técnicas de \glsentrytext{dl} ha representado un gran avance. Sin embargo, estos modelos enfrentan desafíos significativos, principalmente el elevado coste computacional y de memoria.

La teoría cuántica ofrece nuevos paradigmas para el modelado matemático y la computación. Se han desarrollado modelos inspirados en la mecánica cuántica para abordar tareas complejas como la recuperación de información y la desambiguación semántica.

Adicionalmente, en el ámbito del \glsentrytext{ml}, se ha demostrado teóricamente que los modelos cuánticos de \glsentrytext{ml} (\glsentrytext{qml}) presentan una capacidad de expresividad superior a la de los modelos clásicos en ciertos regímenes, haciendo que sea un foco de estudio actual para la mejora de los modelos \glsentrytext{nlp}. El \glsentrytext{qnlp} busca explotar ventajas cuánticas como el paralelismo y la alta dimensionalidad del espacio de estados para procesar lenguaje de manera más eficiente o con mayor riqueza semántica.

\section{Paralelismo entre Espacio Semántico y Espacio de Hilbert}

\section{Codificación de información clásica en estados cuánticos}

\subsection{Basis Encoding y Amplitude Encoding}

\section{Circuitos Cuánticos Variacionales (\glsentrytext{vqc}) como modelos de aprendizaje}